%% ma: L11-B09-0012
%% lop: 11
%% chuong: 3
%% bai: 9
%% dang: 11.9.4
%% loai: TLN
%% muc: TH
%% dap_an: "7,3"
%% nguon: "de-on-tap-giua-ki-1-toan-11-nam-2026-2027-truong-thpt-ngo-quyen-da-nang.pdf (M09, Đề ôn tập giữa kì 1 Toán 11 THPT Ngô Quyền, Đà Nẵng), Mã đề 0112, Phần III Câu 4"
%% kien_thuc: "Bài 9 (các số đặc trưng đo xu thế trung tâm)"
%% trang_thai: chinh-thuc
\begin{ex}
Điểm kiểm tra môn Toán của một nhóm học sinh được cho bởi bảng số liệu ghép nhóm sau:
\begin{center}\small\begin{tabular}{|l|c|c|c|c|c|c|}\hline
Nhóm điểm & $[4;5)$ & $[5;6)$ & $[6;7)$ & $[7;8)$ & $[8;9)$ & $[9;10]$\\\hline
Số học sinh & $3$ & $5$ & $12$ & $15$ & $10$ & $5$\\\hline
\end{tabular}\end{center}
Tính điểm trung bình của nhóm học sinh này (làm tròn đến hàng phần mười).
\shortans{7,3}
\loigiai{$n=50$. \[\bar x=\dfrac{4{,}5\cdot3+5{,}5\cdot5+6{,}5\cdot12+7{,}5\cdot15+8{,}5\cdot10+9{,}5\cdot5}{50}=\dfrac{364}{50}=7{,}28\approx7{,}3.\]}
\end{ex}
